\documentclass{article}
\usepackage{amsmath}
\usepackage{amssymb}
\usepackage{amsthm}
\usepackage{mathtools}

\usepackage{color}
\newcommand{\red}[1]{ {\color{red}{#1}} }
\newcommand{\blue}[1]{ {\color{blue}{#1}} }

\newcommand{\N}{\mathbb{N}}
\newcommand{\Z}{\mathbb{Z}}
\newcommand{\Q}{\mathbb{Q}}
\newcommand{\R}{\mathbb{R}}
\newcommand{\C}{\mathbb{C}}

\theoremstyle{plain}
    \newtheorem{definition}{Definition}
    \newtheorem{example}{Example}
    \newtheorem{proposition}{Proposition}

\title{Writeup}
\author{Caleb and Jeremy}
\date{September 2026}

\begin{document}
\maketitle

\section{}


\subsection{Introduction}

The goal of this section is to characterize the conditions under which the function defined by
\[ 
    \Z/n\Z \xrightarrow[x \mapsto ax+b]{} Z/n\Z
\]
is invertible.
That will give conditions for an \textit{affine cipher} to be a valid encryption method.

To do this we will need some basic number theory, and the following correspondence
between English letters and integers:
\begin{center}
\begin{tabular}{c|c|c|c|c|c}
    A & B & C & $\cdots$ & Y & Z  \\ \hline
    0 & 1 & 2 & $\cdots$ & 24 & 25
\end{tabular}
\end{center}

To motivate our study, we first prove the following fact.

\begin{proposition}
    Let $a,b \in \Z$ and consider the function $f: \Z \to \Z$ defined by 
    \[
        f(x) \coloneqq ax + b.
    \]
    Then $f$ is invertible if and only if $a=\pm1$.
\end{proposition}

\blue{
\begin{proof}
    Suppose $g$ is the inverse of $f$ and $g(x) = cx + d$.....
\end{proof}
}

What is special about $a=\pm1$?
In this case, $a$ is what we call a \textbf{unit} of the ring $\Z$.
We will work to characterize the units of the ring $\Z/n\Z$.
To accomplish this we'll need to build a few ideas of number theory:
\begin{itemize}
    \item Greatest common divisor 
    \item Relative Primality
\end{itemize}




\subsection{Residues}

The symbol $\Z$ denotes the set
\[
    \{ \dots, -2, -1, 0, 1, 2, \dots\}
\]
of \textbf{integers} or \textbf{whole numbers}.

\begin{definition}
    Let $n$ be a positive integer.
    Define $\Z/n\Z$ to be the set 
    \[
        \{ [0], [1], \dots, [n-1] \}.
    \]
    Sometimes, instead of writing $[k]$ we will write $k\pmod{n}$
    or simply $k$ if there is no confusion.
\end{definition}

In some sense, we can think of $\Z/n\Z$ as a collection of all possible 
remainders upon division by $n$.
This idea is illustrated in the following fact.

\begin{proposition}[Division]\label{prop:div}
    Let $a,n \in \Z$ be two positive integers.
    There is exactly one pair $q,r\in\Z$ of positive integers satisfying 
    \begin{enumerate}
        \item $a = qn+r$,
        \item $0\leq r < n$.
    \end{enumerate}
\end{proposition}

\blue{
\begin{proof}
    To do.
\end{proof}
}

In the case when $r=0$, we say $n|a$, pronounces \textbf{$n$ divides $a$}.\footnote{
    The symbol $|$ has an alternative pronunciation: ``guzinta.''
As in, ``$n$ guzinta $a$,'' or, ``$n$ goes into $a$.''}
We refer to $r$ as the remainder upon division by $n$, or 
the \textbf{residue of $a$ modulo $n$}.
We express this by writing $a \equiv r \pmod{n}$.

Two integers may have the same residue modulo $n$.
If $a \equiv r \pmod{n}$, then $a+n \equiv r \pmod{n}$ since 
\[ 
    a+n = (q+1)n + r.
\] 
A similar calculation shows that
\[
    r \equiv a+2n \equiv a+3n \equiv \cdots \pmod{n}.
\]
Thus we think of $r$ as ``representing'' the family $a+kn$ of integers, where $k$
ranges across all elements of $\Z$.



\begin{definition}[Algebraic operations of $\Z/n\Z$]
    Let $[a], [b] \in \Z/n\Z$.
    By Proposition~\ref{prop:div} we have the following two expressions:
    \begin{enumerate}
        \item $a+b = q_1 n + r_1$
        \item $ab = q_2 n + r_2$
    \end{enumerate}
    where $0\leq r_1,r_2 < n$.
    We define an additive and multiplicative operations on $\Z/n\Z$ by 
    \[  
        [a]+[b] \coloneq [r_1], \quad [a] \cdot [b] \coloneq [r_2].
    \]
\end{definition}

These algebraic operations are much nicer than their definition might first suggest.

\begin{proposition}[$\Z/n\Z$ is a ring]
    Let $[a],[b],[c] \in \Z/n\Z$:
    \begin{enumerate}
        \item $[a]+[b] = [b]+[a]$
        \item $([a]+[b]) + [c] = [a] + ([b]+[c])$
        \item $[0]+[a] = [a] = [a]+[0]$
        \item $\forall [a], \exists [b]; [a]+[b] = [0]$
    \end{enumerate}

    \begin{enumerate}
        \item $[a]([b]+[c]) = [a][b] + [a][c]$
        \item $([a]+[b])[c] = [a][c] + [b][c]$
    \end{enumerate}
\end{proposition}

\blue{
\begin{proof}
    To do.
\end{proof}
}




\subsection{Primes}

Now we return to the idea of divisibility.
\begin{definition}
    Let $a,b \in \Z$ be positive integers.
    If $a|b$ we say $b$ is a \textbf{multiple} of $a$, and that 
    $a$ is a \textbf{divisor} of $b$.
    Note that, for all integers $b$, we have $1|b$ and $b|b$.
    A positive integer $p$ is called \textbf{prime} if $p$ has exactly 
    two divisors: $1$ and $p$.
    A non-prime integer is called \textbf{composite}.
\end{definition}


















\subsection{Caesar Ciphers and Affine Ciphers}

We are now in a position to state the result alluded to at the outset.

\begin{proposition}
    The function $f: \Z/n\Z \to \Z/n\Z$ defined by 
    \[
        f(x) \coloneq ax + b
    \]
    is invertible if and only if $\gcd(a,n)=1$.
\end{proposition}

\blue{
\begin{proof}
    Todo.
\end{proof}
}



\end{document}